% MASTER ARM ELECTRONICS. The pin assignment is the figure's real content,
% because the A4/A5 exclusion is what makes 16 analog channels sufficient.
\begin{tikzpicture}[x=1mm, y=1mm, font=\scriptsize]

\node[draw, thick, fill=gray!12, minimum width=44mm, minimum height=62mm,
      align=center] (mcu) at (0,0)
  {\textbf{Teensy 4.1}\\[2pt] i.MX RT1062\\ \SI{600}{\mega\hertz}\\[4pt]
   ADC \SI{12}{bit}\\ \texttt{analogReadAveraging(1)}\\[4pt]
   loop \SI{10}{\milli\second}\\ USB serial\\ \SI{115200}{baud}};

% ---- left: K1 pots
\node[draw, fill=blue!8, minimum width=30mm, align=center] (k1) at (-52, 20)
  {\textbf{K1 (left) pots}\\ 7 channels\\ \texttt{A0 A1 A2 A3}\\ \texttt{A6 A7 A8}};
\node[draw, fill=blue!8, minimum width=30mm, align=center] (k2) at (-52, 0)
  {\textbf{K2 (right) pots}\\ 7 channels\\ \texttt{A9 A10 A11}\\ \texttt{A12 A13 A16 A17}};
\node[draw, fill=blue!8, minimum width=30mm, align=center] (fsr) at (-52, -18)
  {\textbf{FSRs}\\ \texttt{A14} (left)\\ \texttt{A15} (right)};
\node[draw, fill=blue!8, minimum width=30mm, align=center] (btnn) at (-52, -34)
  {\textbf{buttons}\\ digital \texttt{2}, \texttt{3}\\ \texttt{INPUT\_PULLUP}};

% ---- right: I2C
\node[draw, fill=red!8, draw=red!60!black, thick, minimum width=32mm,
      align=center] (wire) at (52, 14)
  {\textbf{Wire (I\textsuperscript{2}C0)}\\ SDA pin 18 $=$ \texttt{A4}\\
   SCL pin 19 $=$ \texttt{A5}\\ \SI{400}{\kilo\hertz}};
\node[draw, fill=blue!8, minimum width=32mm, align=center] (imu1) at (52, -4)
  {\textbf{MPU-6050 \#1}\\ \texttt{0x68} (AD0 low)\\ left wrist};
\node[draw, fill=blue!8, minimum width=32mm, align=center] (imu2) at (52, -20)
  {\textbf{MPU-6050 \#2}\\ \texttt{0x69} (AD0 high)\\ right wrist};
\node[draw, fill=green!8, draw=green!45!black, minimum width=32mm,
      align=center] (host) at (52, -38)
  {\textbf{host}\\ \texttt{/dev/ttyACM*}\\ auto-detected};

\draw[flow] (k1.east)   -- (mcu.west |- k1.east);
\draw[flow] (k2.east)   -- (mcu.west |- k2.east);
\draw[flow] (fsr.east)  -- (mcu.west |- fsr.east);
\draw[flow] (btnn.east) -- (mcu.west |- btnn.east);
\draw[flow] (mcu.east |- wire.west) -- (wire.west);
\draw[thick] (wire.south) -- (imu1.north);
\draw[thick] (imu1.south) -- (imu2.north);
\draw[flow] (mcu.east |- host.west) -- (host.west);

% ---- the exclusion, called out
\node[draw, dashed, thick, red!60!black, fill=red!4, text width=74mm,
      align=left, font=\scriptsize] at (0, -50)
  {\textbf{\color{red!60!black}Why \texttt{A4} and \texttt{A5} carry no
   potentiometer.} The default \texttt{Wire} bus is pins 18/19, which
   \emph{are} \texttt{A4}/\texttt{A5}. An earlier revision listed both inside
   \texttt{potPinsK1}, so two pot channels shared conductors with the IMU data
   bus. Teensy~4.1 exposes 18 analog-capable pins; removing two leaves exactly
   16, which is what $14 + 2$ requires. \texttt{Wire1} (16/17) and
   \texttt{Wire2} (24/25) are never initialised, so \texttt{A2}, \texttt{A3},
   \texttt{A10} and \texttt{A11} are free as plain analog inputs.};

\end{tikzpicture}
